\label{sec:method}

\subsection{方法概述}

现有刚体接触仿真中，基于少量接触点缓存的离散接触流形更强调时间连续性与求解效率，而非连续接触区域的高保真力学表征。与之相对，基于接触曲面与压力场积分的思路体现了“surface-first”的建模思想：先构造接触面，再由接触面上的分布量得到合力与合力矩。进一步地，等效接触点（Equivalent Contact Point, ECP）类方法表明，对单个凸接触区域，可以用一个等效接触点及其对应的接触 wrench 来统一线接触与面接触。上述认识共同说明：高精度刚体接触的关键，不在于缓存少量几何点，而在于如何在刚体动力学层面保留连续接触区域的等效 wrench。

基于此，本文提出一种面向 Chrono NSC 的\textbf{多 Patch / 多 Subpatch 保持 Wrench 等效的 SDF 接触约化方法}。该方法以 SDF 为几何基础，首先为每个候选接触对构造高密度接触区域参考模型；随后将连续接触区域分解为多个 patch 与 subpatch；最后对每个 subpatch 构造少量代表点组成的 reduced support set，使其在给定误差阈值下逼近 dense 接触模型所张成的可行 wrench 锥，并将所得代表接触点重新注入 Chrono NSC 接触容器。该方法定位于碰撞--接触建模接口层，而不修改全局 NSC 互补求解框架。

本文方法的核心思想可概括为：
\begin{enumerate}[label=(\arabic*)]
    \item \textbf{surface-first}：先由 SDF 构造高密度接触区域，而非直接信任离散点流形；
    \item \textbf{wrench-first}：约化目标不是保留若干几何点，而是逼近 dense 接触模型的可行 wrench 锥；
    \item \textbf{adaptive refinement}：以几何误差、wrench 误差、CoP 误差与连续区域 gap 误差共同驱动自适应细分；
    \item \textbf{solver-compatible}：最终输出仍为普通点接触，从而直接复用 Chrono NSC 的求解链路。
\end{enumerate}

上述四点刻画了本文的瞬时约化目标。然而，滑移主导场景的数值现象表明：仅保持某一时刻的合力、合力矩与 CoP 仍不足以保证下一时间步的速度更新一致。为此，本文在瞬时 wrench-preserving 框架之上，进一步引入\textbf{一步动力学一致性（one-step dynamic consistency）}：dense 接触、reduced support 以及 solver-facing reinjection 三者不仅要在瞬时 wrench 上保持等效，还应在目标 NSC 时间步推进下产生近似一致的广义冲量与广义速度更新。

\subsection{问题定义}

考虑两个刚体 $A$ 与 $B$，其几何分别由符号距离场 $\phi_A(x)$ 与 $\phi_B(x)$ 描述。对于任一候选接触对，在相对位姿给定时构造局部 gap 场 $g(x)$。在本文中，$g(x)$ 可由双 SDF 组合、单边表面采样或其他等价构造获得，其目标是给出局部非穿透约束的几何刻画。接触区域定义为
\begin{equation}
\Omega = \{x \in \mathbb{R}^3 \mid g(x) \le 0\}.
\end{equation}

为获得高精度参考接触模型，将 $\Omega$ 离散为带面积权重的高密积分点集合
\begin{equation}
\mathcal{Q} = \{(x_k, a_k, n_k)\}_{k=1}^{N_q},
\end{equation}
其中 $x_k$ 为接触积分点，$a_k$ 为面积权重，$n_k$ 为由 SDF 梯度给出的局部单位法向，即
\begin{equation}
 n_k = \frac{\nabla \phi(x_k)}{\|\nabla \phi(x_k)\|}.
\end{equation}

在每个积分点上，引入局部接触力变量
\begin{equation}
\lambda_k =
\begin{bmatrix}
 t_{k1} \\
 t_{k2} \\
 n_k^c
\end{bmatrix},
\qquad
n_k^c \ge 0,
\qquad
\sqrt{t_{k1}^2 + t_{k2}^2} \le \mu_k n_k^c,
\end{equation}
其中 $n_k^c$ 表示法向反力标量，$(t_{k1}, t_{k2})$ 表示局部切向反力分量，$\mu_k$ 为局部摩擦系数。

对于参考点 $c$，记 $x_k-c$ 在局部切平面基 $(e_1,e_2,e_3)$ 中的坐标为 $(u_k,v_k,0)$。则每个积分点对局部六维 wrench 的映射矩阵可写为
\begin{equation}
G(x_k) =
\begin{bmatrix}
1 & 0 & 0 \\
0 & 1 & 0 \\
0 & 0 & 1 \\
0 & 0 & v_k \\
0 & 0 & -u_k \\
-v_k & u_k & 0
\end{bmatrix}.
\end{equation}

因此，dense 接触模型可产生的任意局部 wrench 为
\begin{equation}
 w = \sum_{k=1}^{N_q} G(x_k)\lambda_k.
\end{equation}

进一步定义 dense 可行 wrench 锥
\begin{equation}
\mathcal{C}_{\mathrm{dense}}
=
\left\{
\sum_{k=1}^{N_q} G(x_k)\lambda_k
\;\middle|\;
 n_k^c \ge 0,\;
 \|(t_{k1},t_{k2})\|_2 \le \mu_k n_k^c
\right\}.
\end{equation}

本文约化的目标不是保留 $\mathcal{Q}$ 的全部几何点，而是构造一个低维 reduced support set，使其对应的 reduced wrench 锥尽可能逼近 $\mathcal{C}_{\mathrm{dense}}$。

\subsection{多 Patch / 多 Subpatch 分解}

\subsubsection{Patch 分解}

由于一个接触对可能包含多个不连通接触区域，本文首先在高密积分点集合 $\mathcal{Q}$ 上建立邻接图
\begin{equation}
(i,j)\in E
\iff
\|x_i - x_j\| < h
\;\land\;
n_i \cdot n_j > \cos\theta_0
\;\land\;
g(x_i)\le 0,\; g(x_j)\le 0,
\end{equation}
其中 $h$ 为局部采样尺度，$\theta_0$ 为法向相容阈值。图的连通分量定义为接触 patch：
\begin{equation}
\Omega = \bigcup_{\ell=1}^{N_p} \Omega_\ell.
\end{equation}

这种处理确保不同的连通接触区域不会被错误合并为单一 patch。

\subsubsection{Subpatch 分解}

仅按连通性分块仍不足以保证局部 wrench 等效的精度。为此，本文进一步将每个 patch $\Omega_\ell$ 分解为若干 subpatch：
\begin{equation}
\Omega_\ell = \bigcup_{r=1}^{m_\ell} \Omega_{\ell r}.
\end{equation}

subpatch 的细分由如下误差指标联合驱动：
\begin{enumerate}[label=(\roman*)]
    \item \textbf{局部平面化误差}
    \begin{equation}
    e_{\mathrm{plane}} = \kappa_{\max} D,
    \end{equation}
    其中 $D$ 为当前块的直径，$\kappa_{\max}$ 为最大主曲率；

    \item \textbf{二阶矩误差}
    \begin{equation}
    e_{\Sigma}
    =
    \frac{\|\Sigma_{\mathrm{red}}-\Sigma_{\mathrm{dense}}\|_F}
    {\|\Sigma_{\mathrm{dense}}\|_F+\varepsilon};
    \end{equation}

    \item \textbf{可行 wrench 锥误差}
    \begin{equation}
    e_{\mathcal{C}}
    =
    \max_{s\in\mathcal{S}_w}
    \frac{|h_{\mathrm{dense}}(s)-h_{\mathrm{red}}(s)|}
    {h_{\mathrm{dense}}(s)+\varepsilon},
    \end{equation}
    其中 $h_{\mathcal{C}}(s)=\sup_{w\in\mathcal{C}} s^\top w$ 为 support function；

    \item \textbf{内部漏穿透误差}
    \begin{equation}
    e_g = \max_{s_i\in \mathcal{S}_{\ell r}} \max(0,-g(s_i)),
    \end{equation}
    其中 $\mathcal{S}_{\ell r}$ 为布置于 subpatch 内部的 sentinel 点集合。
\end{enumerate}

一旦任一误差超过预设阈值，当前块将继续沿主方向、法向变化方向或局部骨架方向进行细分。由此，本文形成由误差控制的 patch-subpatch 自适应层级结构。

\subsection{局部保持 Wrench 等效的 Reduced Support Set}

\subsubsection{局部坐标系与 Reduced Wrench 锥}

对于任意 subpatch $\Omega_{\ell r}$，定义其局部参考点 $c_{\ell r}$ 为面积加权质心，局部法向 $e_3$ 为面积加权平均法向，$e_1,e_2$ 为张成局部切平面的正交基。设 reduced support set 为
\begin{equation}
Q_{\ell r}=\{q_j=(u_j,v_j,0)\}_{j=1}^{m},
\end{equation}
其中 $m$ 为当前 subpatch 的代表点数。对应的 reduced 点力变量为
\begin{equation}
\hat{\lambda}_j=
\begin{bmatrix}
\hat{t}_{j1}\\
\hat{t}_{j2}\\
\hat{n}_j
\end{bmatrix},
\qquad
\hat{n}_j\ge0,
\qquad
\sqrt{\hat{t}_{j1}^2+\hat{t}_{j2}^2}\le \mu_j \hat{n}_j.
\end{equation}

则 reduced 模型对应的可行 wrench 锥为
\begin{equation}
\mathcal{C}_{\mathrm{red}}(Q_{\ell r})
=
\left\{
\sum_{j=1}^{m} G(q_j)\hat{\lambda}_j
\right\}.
\end{equation}

我们的目标是使 $\mathcal{C}_{\mathrm{red}}$ 在给定容忍度下逼近 $\mathcal{C}_{\mathrm{dense}}$。

\subsubsection{三点局部 Wrench 完备性}

当 $\Omega_{\ell r}$ 局部近似共面时，三个非共线代表点已足以生成一个完整的局部六维 wrench 基。具体地，总 wrench 分量满足
\begin{equation}
\begin{aligned}
T_x &= \sum_{j=1}^{m} \hat{t}_{j1},\\
T_y &= \sum_{j=1}^{m} \hat{t}_{j2},\\
N   &= \sum_{j=1}^{m} \hat{n}_j,\\
M_x &= \sum_{j=1}^{m} v_j \hat{n}_j,\\
M_y &= -\sum_{j=1}^{m} u_j \hat{n}_j,\\
M_z &= \sum_{j=1}^{m}\left(u_j \hat{t}_{j2}-v_j \hat{t}_{j1}\right).
\end{aligned}
\end{equation}

对 $m=3$ 且三点非共线的情形，法向分量满足
\begin{equation}
\begin{bmatrix}
1 & 1 & 1\\
v_1 & v_2 & v_3\\
-u_1 & -u_2 & -u_3
\end{bmatrix}
\begin{bmatrix}
\hat{n}_1\\
\hat{n}_2\\
\hat{n}_3
\end{bmatrix}
=
\begin{bmatrix}
N\\
M_x\\
M_y
\end{bmatrix},
\end{equation}
切向分量满足
\begin{equation}
\begin{bmatrix}
1 & 0 & 1 & 0 & 1 & 0\\
0 & 1 & 0 & 1 & 0 & 1\\
-v_1 & u_1 & -v_2 & u_2 & -v_3 & u_3
\end{bmatrix}
\begin{bmatrix}
\hat{t}_{11}\\
\hat{t}_{12}\\
\hat{t}_{21}\\
\hat{t}_{22}\\
\hat{t}_{31}\\
\hat{t}_{32}
\end{bmatrix}
=
\begin{bmatrix}
T_x\\
T_y\\
M_z
\end{bmatrix}.
\end{equation}

当三点非共线时，上述两个矩阵均满秩，因此在忽略库仑锥约束的情况下，三点已对局部六维 wrench 完备。由此可知，\textbf{三点是局部共面 subpatch 的最小 wrench 完备基}。若摩擦可行性不足或几何条件恶化，则提高 $m$ 或继续细分 subpatch。

\subsubsection{法向可行性与 CoP 支撑域}

对 $N>0$，定义中心压力点（center of pressure, CoP）
\begin{equation}
 u_c = -\frac{M_y}{N},
 \qquad
 v_c = \frac{M_x}{N}.
\end{equation}

令 $\lambda_j = \hat{n}_j/N$，则有
\begin{equation}
\lambda_1+\lambda_2+\lambda_3=1,
\qquad
\sum_{j=1}^{3}\lambda_j q_j = (u_c,v_c,0).
\end{equation}

\begin{proposition}
对于三点 reduced support set，存在 $\hat{n}_j\ge 0$ 使法向合力与法向力矩精确匹配，当且仅当 CoP 落在三角形 $\mathrm{conv}\{q_1,q_2,q_3\}$ 内。
\end{proposition}

该命题给出了一个重要设计原则：代表点三角形并非任意选取，而必须覆盖 dense subpatch 在给定工况下可能出现的 CoP 活动区域。

\subsection{代表点位置优化}

\subsubsection{几何初始化}

首先计算 dense subpatch 投影到局部平面的二维点集 $\xi_k$ 的一、二阶矩：
\begin{equation}
A = \sum_k a_k,
\qquad
\bar{\xi} = \frac{1}{A}\sum_k a_k \xi_k,
\qquad
\Sigma = \frac{1}{A}\sum_k a_k(\xi_k-\bar{\xi})(\xi_k-\bar{\xi})^\top.
\end{equation}

令 $LL^\top=\Sigma$。对初始三点，可取标准二维 simplex 顶点 $s_j$ 构造
\begin{equation}
 q_j^{(0)} = \bar{\xi} + \gamma L s_j,
\end{equation}
其中 $\gamma$ 为尺度因子。该初始化使 reduced support set 在均值与二阶尺度上与 dense patch 对齐。

\subsubsection{精度优先目标函数}

最终代表点通过如下优化求得：
\begin{equation}
\min_{Q\subset \mathcal{H}}
\;
\alpha \|c(Q)-\bar{\xi}\|^2
+\beta \|\Sigma(Q)-\Sigma\|_F^2
+\gamma E_{\mathcal{C}}(Q)^2
-\eta \log\det(B_n B_n^\top)
-\zeta \log\det(B_t B_t^\top),
\end{equation}
其中：
\begin{itemize}[leftmargin=2em]
    \item $\mathcal{H}$ 为当前 subpatch 的几何支撑域；
    \item $c(Q)$ 与 $\Sigma(Q)$ 分别为代表点集合的一、二阶矩；
    \item $E_{\mathcal{C}}(Q)$ 为 dense/reduced wrench 锥的 support-function 误差；
    \item $B_n$ 与 $B_t$ 分别为法向与切向矩阵。
\end{itemize}

其中前两项约束 reduced support set 的几何尺度，第三项直接最小化 wrench 锥误差，后两项提升法向与切向矩阵的条件数，避免代表点塌缩或杠杆臂退化。

\subsection{局部参考 Wrench 预测与 Reduced 力分配}

\subsubsection{局部参考 Wrench 预测}

仅逼近可行 wrench 锥仍不足以决定当前时间步应落在锥中的哪一点。因此，本文为每个 dense subpatch 构造一个\textbf{局部参考 wrench 预测器}，得到当前时刻参考 wrench $w_{\ell r}^\star$。

设当前相对 twist 为 $(v_{\mathrm{rel}}, \omega_{\mathrm{rel}})$。本文采用局部 dense 微求解方式，在 dense 积分点上解一个小型 SOCP/QP：
\begin{equation}
\min_{\lambda}
\frac{1}{2}\lambda^\top H \lambda + g^\top \lambda
\quad
\text{s.t.}\quad
\lambda \in \mathcal{K}_\mu,
\end{equation}
其中 $\mathcal{K}_\mu$ 表示由所有积分点局部库仑锥组成的直积锥，$H$ 与 $g$ 来自局部相对速度、局部线性化几何与近场接触响应。解得 $\lambda^\star$ 后，定义参考 wrench
\begin{equation}
 w_{\ell r}^\star
 =
 \sum_{k} G(x_k)\lambda_k^\star.
\end{equation}

这里的局部 dense 微求解并不替代全局 NSC，而只是为 reduced support set 提供当前时刻的高精参考目标。

\subsubsection{Reduced 力分配}

给定优化后的代表点集合 $Q_{\ell r}$ 与参考 wrench $w_{\ell r}^\star$，求解 reduced 点力：
\begin{equation}
\min_{\hat{\lambda}}
\frac{1}{2}\hat{\lambda}^\top W \hat{\lambda}
\end{equation}
\begin{equation}
\text{s.t.}\quad
\sum_{j=1}^{m} G(q_j)\hat{\lambda}_j = w_{\ell r}^\star,
\end{equation}
\begin{equation}
\hat{n}_j \ge 0,
\qquad
\sqrt{\hat{t}_{j1}^2+\hat{t}_{j2}^2}\le \mu_j \hat{n}_j.
\end{equation}

该问题同样是一个小型 SOCP。若可行，则说明当前 reduced support set 能对该 subpatch 的当前参考 wrench 实现\textbf{精确等效}；若不可行，则说明当前基不够强，必须执行以下操作之一：
\begin{enumerate}[label=(\arabic*)]
    \item 增加代表点数 $m$；
    \item 扩大代表点支撑域并重新优化；
    \item 将当前 subpatch 继续分裂。
\end{enumerate}

本文不允许在 reduced 力分配阶段以牺牲精度的方式“硬拟合”不可行解，而是始终以误差控制驱动模型升阶或细分。

\subsection{连续区域内部非穿透监测}

少量代表点接触方法的一个典型问题是：虽然代表点满足局部非穿透，但 patch 内部可能仍发生漏穿透。为此，本文在每个 subpatch 内布置一组不进入求解器的 sentinel 点
\begin{equation}
 \mathcal{S}_{\ell r} = \{s_i\}_{i=1}^{N_s},
\end{equation}
并在每一步预测后计算
\begin{equation}
 e_g
 =
 \max_{s_i\in\mathcal{S}_{\ell r}}
 \max(0,-g(s_i)).
\end{equation}

若 $e_g$ 超过阈值 $\varepsilon_g$，则说明 reduced support set 虽然在离散支撑点上满足约束，但未能充分覆盖连续接触区域的内部；此时优先对当前 subpatch 继续细分，并在必要时增加 support points。

\subsection{时间连续性与 Warm Start}

由于接触区域在相邻时间步之间通常具有时序连续性，本文为每个 subpatch 维护 persistent identifier，并在相邻步之间进行匹配：
\begin{equation}
\mathrm{match}(\Omega_{\ell r}^{t}, \Omega_{\ell' r'}^{t+\Delta t}).
\end{equation}

匹配指标包括质心距离、平均法向夹角、面积重叠率、局部骨架重叠与代表点集合的 Hausdorff 距离。对匹配成功的 subpatch，将上一时间步的 reduced impulses $\hat{\lambda}^{\mathrm{old}}$ 传输至新的 support set，传输过程同样采用保持 wrench 等效的优化：
\begin{equation}
\min_{\hat{\lambda}^{\mathrm{new}}}
\|\hat{\lambda}^{\mathrm{new}}\|_W^2
\quad
\text{s.t.}\quad
\sum_j G(q_j^{\mathrm{new}})\hat{\lambda}_j^{\mathrm{new}}
=
\sum_j G(q_j^{\mathrm{old}})\hat{\lambda}_j^{\mathrm{old}}.
\end{equation}

这意味着 warm start 的传递对象不是逐点冲量值本身，而是对应的等效 wrench。

\subsection{Dense 一步冲量模型}

前述 dense 微求解给出了瞬时参考 wrench。为了把该参考量直接连接到目标求解器的一步时间推进，本文进一步在每个 dense subpatch 上定义一步冲量模型。设局部广义速度为
\begin{equation}
\nu =
\begin{bmatrix}
v \\
\omega
\end{bmatrix},
\end{equation}
时间步长为 $h$，并对每个 dense 积分点采用摩擦锥的多射线近似。记第 $k$ 个积分点的射线权重为 $\alpha_k \ge 0$，射线方向矩阵为 $R_k$，则一步接触冲量可写为
\begin{equation}
\Delta p_{\mathrm{dense}}
=
\sum_{k=1}^{N_q} G(x_k) \, h R_k \alpha_k
=
A_{\mathrm{dense}}\alpha.
\end{equation}

令 $J_{\mathrm{dense}}$ 为 dense 积分点的局部接触 Jacobian，$M$ 为刚体系统在当前构型下的广义质量矩阵，则一步局部相对速度的线性化更新满足
\begin{equation}
u^{+}
\approx
u^{-} + B_{\mathrm{dense}}\alpha,
\qquad
B_{\mathrm{dense}} = J_{\mathrm{dense}} M^{-1} J_{\mathrm{dense}}^{\top} R_{\mathrm{dense}}.
\end{equation}
同理，对局部 gap 预测写成
\begin{equation}
d^{+}
\approx
d^{-} + C_{\mathrm{dense}}\alpha,
\end{equation}
其中 $C_{\mathrm{dense}}$ 由法向运动预测与局部几何线性化给出。

在摩擦锥多射线化与单边 gap 的活跃集线性化下，dense 一步微求解可统一写为
\begin{equation}
\min_{\alpha \ge 0}
\;
\frac{1}{2}\alpha^{\top} H_{\mathrm{dense}}\alpha
+
g_{\mathrm{dense}}^{\top}\alpha,
\label{eq:dense-micro-qp}
\end{equation}
其中
\begin{equation}
H_{\mathrm{dense}}
=
B_{\mathrm{dense}}^{\top}W_u B_{\mathrm{dense}}
+
C_{\mathrm{dense}}^{\top}W_g C_{\mathrm{dense}}
+
\rho_{\alpha} I,
\end{equation}
\begin{equation}
g_{\mathrm{dense}}
=
B_{\mathrm{dense}}^{\top}W_u u^{-}
+
C_{\mathrm{dense}}^{\top}W_g d^{-}
-
\rho_{\alpha}\alpha_{\mathrm{tr}},
\end{equation}
$W_u$ 与 $W_g$ 分别衡量局部速度与 gap 违反的代价，$\alpha_{\mathrm{tr}}$ 为上一时间步传输而来的局部冲量初值。求得最优解 $\alpha^\star$ 后，定义 dense 参考一步冲量为
\begin{equation}
\Delta p_{\ell r}^{\star}
=
A_{\mathrm{dense}}\alpha^\star.
\end{equation}

与前述瞬时参考 wrench 不同，$\Delta p_{\ell r}^{\star}$ 直接对应当前时间步中 dense subpatch 对刚体速度更新的最优局部冲量解释，因此它是后续 reduced solve 与 reinjection 的首要目标量。

\subsection{Reduced 冲量集近似}

对任意 subpatch $\Omega_{\ell r}$ 及其 reduced support set $Q_{\ell r}$，同样在每个 support 上采用摩擦锥射线近似。记 reduced 射线权重为 $\beta \ge 0$，则 reduced 一步冲量写为
\begin{equation}
\Delta p_{\mathrm{red}}
=
A_{\mathrm{red}}(Q_{\ell r})\beta.
\end{equation}

由此可定义 dense 与 reduced 的一步可行冲量集：
\begin{equation}
\mathcal{I}_{\mathrm{dense}}^{h}
=
\left\{
A_{\mathrm{dense}}\alpha
\;\middle|\;
\alpha \ge 0
\right\},
\qquad
\mathcal{I}_{\mathrm{red}}^{h}(Q_{\ell r})
=
\left\{
A_{\mathrm{red}}(Q_{\ell r})\beta
\;\middle|\;
\beta \ge 0
\right\}.
\end{equation}

于是，本文不再仅要求 $\mathcal{C}_{\mathrm{red}}$ 逼近 $\mathcal{C}_{\mathrm{dense}}$，还要求在一步时间尺度上，$\mathcal{I}_{\mathrm{red}}^{h}$ 对 $\mathcal{I}_{\mathrm{dense}}^{h}$ 的主要支撑方向保持一致。定义冲量支持函数
\begin{equation}
h_{\mathcal{I}}(s)
=
\sup_{\Delta p \in \mathcal{I}} s^{\top}\Delta p,
\end{equation}
则冲量锥误差写为
\begin{equation}
e_{\mathcal{I}}
=
\max_{s\in\mathcal{S}_{I}}
\frac{|h_{\mathrm{dense}}^{I}(s)-h_{\mathrm{red}}^{I}(s)|}
{h_{\mathrm{dense}}^{I}(s)+\varepsilon}.
\end{equation}

对应的 reduced 分配问题为
\begin{equation}
\min_{\beta \in \mathcal{K}_{\mu}}
\;
\frac{1}{2}
\|W_p(A_{\mathrm{red}}\beta-\Delta p_{\ell r}^{\star})\|_2^2
+
\frac{\rho_{\beta}}{2}\|\beta-\beta_{\mathrm{tr}}\|_2^2,
\label{eq:reduced-impulse-qp}
\end{equation}
其中 $\beta_{\mathrm{tr}}$ 为由上一时间步传输而来的冲量种子。若采用高分辨率摩擦射线，则式~\eqref{eq:reduced-impulse-qp} 可看作对连续摩擦锥 SOCP 的多面体逼近；若保留圆锥约束，则其原始形式就是一个小规模 SOCP。该层的目标不再只是“拟合参考 wrench”，而是显式逼近 dense subpatch 的一步冲量可行域。

\subsection{Solver-Facing Reinjection 优化}

对目标 NSC 求解器而言，reduced support 并不能直接参与求解，而必须进一步转化为离散点接触集合。为此，本文显式定义 reinjection 映射
\begin{equation}
\mathcal{R}:
\left(Q_{\ell r}, \beta\right)
\mapsto
\mathcal{Y}_{\ell r}
=
\{(y_s,n_s,\gamma_s)\}_{s=1}^{N_y},
\end{equation}
其中 $y_s$ 为 emitted slot 位置，$n_s$ 为 slot 法向，$\gamma_s$ 为 slot 冲量变量。

与把 reinjection 当作简单几何模板不同，本文将其视作一个受约束的小规模优化问题：
\begin{equation}
\min_{\{y_s,n_s,\gamma_s\}}
\;
\|W_I(\Delta p_{\mathrm{emit}}-\Delta p_{\mathrm{red}})\|_2^2
+
\|W_T(\tau_{\mathrm{emit}}-\tau_{\mathrm{red}})\|_2^2
+
\chi E_g(\mathcal{Y}_{\ell r})
+
\upsilon E_a(\mathcal{Y}_{\ell r}),
\label{eq:reinjection-opt}
\end{equation}
其中
\begin{equation}
\Delta p_{\mathrm{emit}}
=
\sum_{s=1}^{N_y}\Psi(y_s,n_s)\gamma_s,
\end{equation}
$\tau_{\mathrm{emit}}$ 表示 reinjected 点接触在切向一阶与二阶矩意义下的冲量分布量，$E_g$ 用于惩罚 reinjection 后的内部 gap 违反，$E_a$ 用于惩罚不稳定的激活/关闭切换。约束包括
\begin{equation}
y_s \in \mathrm{conv}(Q_{\ell r}) \oplus \mathbb{B}(0,r_h),
\qquad
n_s \cdot \bar{n}_{\ell r} \ge \cos\theta_h,
\qquad
\gamma_s \in \mathcal{K}_{\mu}^{\mathrm{slot}}.
\end{equation}

式~\eqref{eq:reinjection-opt} 的意义在于：solver-facing reinjection 本身就是一个需要显式控制切向可行域损失的优化层，而不是 reduced support 到点接触的无损映射。它为后文的动态误差分析提供了独立的误差来源。

\subsection{一步动力学一致性}

设 dense 模型、reduced 模型以及 reinjected NSC 点接触在当前时间步得到的广义冲量分别为 $\Delta p_{\mathrm{dense}}^\star$、$\Delta p_{\mathrm{red}}^\star$ 与 $\Delta p_{\mathrm{emit}}^\star$，则其一步速度更新分别满足
\begin{equation}
\nu_{\mathrm{dense}}^{+}
=
\nu^{-} + M^{-1}\Delta p_{\mathrm{dense}}^\star,
\qquad
\nu_{\mathrm{emit}}^{+}
=
\nu^{-} + M^{-1}\Delta p_{\mathrm{emit}}^\star.
\end{equation}

由此可见，本文最终关心的不再只是瞬时 wrench 误差，而是一步广义速度误差
\begin{equation}
e_V
=
\frac{\|\nu_{\mathrm{emit}}^{+}-\nu_{\mathrm{dense}}^{+}\|_2}
{\|\nu_{\mathrm{dense}}^{+}-\nu^{-}\|_2+\varepsilon}.
\end{equation}

\begin{proposition}
若当前构型邻域内的质量矩阵 $M$ 有界且接触 Jacobian 关于状态局部 Lipschitz 连续，则存在常数 $c_M,c_R>0$，使得
\begin{equation}
\|\nu_{\mathrm{emit}}^{+}-\nu_{\mathrm{dense}}^{+}\|_2
\le
c_M
\|\Delta p_{\mathrm{red}}^\star-\Delta p_{\mathrm{dense}}^\star\|_2
+
c_R
\|\Delta p_{\mathrm{emit}}^\star-\Delta p_{\mathrm{red}}^\star\|_2
+
O(h^2).
\end{equation}
\end{proposition}

该命题说明：一步速度更新误差可分解为两层。第一层来自 dense 到 reduced 的冲量近似误差；第二层来自 reduced 到 solver-facing reinjection 的再离散误差。当前实验中最顽固的滑移误差，正对应第二层在切向冲量可行域上的残余差距。

\subsection{冲量一致性误差}

为了把前述理论直接连接到实验与实现，本文在原有 $e_{\mathrm{plane}}$、$e_{\Sigma}$、$e_{\mathcal{C}}$、$e_g$、$e_w$、$e_{\mathrm{CoP}}$ 之外，再引入三类误差：
\begin{equation}
e_I
=
\frac{\|\Delta p_{\mathrm{red}}^\star-\Delta p_{\mathrm{dense}}^\star\|_2}
{\|\Delta p_{\mathrm{dense}}^\star\|_2+\varepsilon},
\qquad
e_V
=
\frac{\|\nu_{\mathrm{emit}}^{+}-\nu_{\mathrm{dense}}^{+}\|_2}
{\|\nu_{\mathrm{dense}}^{+}-\nu^{-}\|_2+\varepsilon},
\end{equation}
\begin{equation}
e_R
=
\max_{s\in\mathcal{S}_{T}}
\frac{|h_{\mathrm{emit}}^{I}(s)-h_{\mathrm{red}}^{I}(s)|}
{h_{\mathrm{red}}^{I}(s)+\varepsilon},
\end{equation}
其中 $e_I$ 衡量 dense 与 reduced 之间的一步冲量误差，$e_V$ 衡量 dense 与 reinjected solver contacts 之间的一步速度更新误差，$e_R$ 衡量 reinjection 在切向冲量可行域上的损失。于是，本文的最终动态目标可表述为
\begin{equation}
\text{static wrench equivalence}
\;+\;
\text{impulse-set preservation}
\;+\;
\text{solver-facing reinjection consistency}.
\end{equation}

\subsection{自适应误差控制}

综合前述定义，本文在每个 subpatch 上维护以下误差：
\begin{equation}
 e_{\mathrm{plane}},\quad
 e_{\mathcal{C}},\quad
 e_{\mathcal{I}},\quad
 e_w=
 \frac{\|w^\star-\hat{w}\|}{\|w^\star\|+\varepsilon},\quad
 e_I,\quad
 e_V,\quad
 e_R,\quad
 e_{\mathrm{CoP}},\quad
 e_g,\quad
 e_{\mathrm{topo}}.
\end{equation}
其中 $e_w$ 衡量参考 wrench 与 reduced wrench 的匹配误差，$e_I$ 与 $e_V$ 分别衡量一步冲量与一步速度更新误差，$e_R$ 衡量 reinjection 的切向可行域缺口，$e_{\mathrm{CoP}}$ 衡量 CoP 的漂移误差。自适应策略如下：
\begin{itemize}[leftmargin=2em]
    \item 若 $e_{\mathrm{plane}}$ 或 $e_{\mathrm{topo}}$ 主导，则细分 subpatch；
    \item 若 $e_{\mathcal{C}}$、$e_{\mathcal{I}}$ 或 $e_{\mathrm{CoP}}$ 主导，则重新优化代表点位置；
    \item 若 reduced 力分配不可行或 $e_I$ 持续偏大，则增加代表点数或提高摩擦射线分辨率；
    \item 若 $e_R$ 主导，则继续优化 solver-facing reinjection 的 slot 位置、法向和切向冲量分配；
    \item 若 $e_g$ 主导，则细分 subpatch 并增加 sentinel 点；
    \item 若 $e_V$ 在多个时间步上持续偏大，则优先调整时间传输与 reinjection，而不是仅继续拟合瞬时 wrench。
\end{itemize}

由此形成如下误差驱动循环：
\begin{enumerate}[label=\arabic*.]
    \item 构造 dense patch；
    \item patch / subpatch 分解；
    \item 代表点初始化与优化；
    \item 局部参考 wrench 预测；
    \item reduced SOCP 力分配；
    \item 误差评估；
    \item 若误差超阈值，则升阶或细分并重复。
\end{enumerate}

该过程体现了本文“\textbf{精度优先}”的设计原则：允许 reduced support set 的规模自适应增加，但不允许在误差失控的条件下维持低维表示。

\subsection{与 Chrono NSC 的集成}

Chrono 的 \verb|ChSystem| 提供了 \verb|SetContactContainer()| 接口，可替换底层 contact container；该 contact container 负责从碰撞检测系统收集接触信息并用于接触力生成。默认 \verb|AddContactCallback| 仅允许在接触被加入容器时修改复合材料参数，而不能替代几何级接触流形重构。因此，本文将所提方法集成为一个自定义 NSC contact container 或自定义接触报告模块，其工作流程为：Bullet 仅负责 broadphase 与碰撞对管理；对每个候选接触对，由本文方法构造 dense SDF 接触区域、执行多 patch / 多 subpatch 约化，并将所得 reduced support points 重新封装为普通 NSC 点接触，最终交由 Chrono NSC 的全局描述器和互补求解器处理。

设最终全部 subpatch 的代表点总数为
\begin{equation}
N_{\mathrm{red}} = \sum_{\ell=1}^{N_p}\sum_{r=1}^{m_\ell} m_{\ell r}.
\end{equation}

则系统级接触自由度规模由原始 dense 模型的 $O(N_q)$ 降至 $O(N_{\mathrm{red}})$，同时由于每个代表点仍以标准点接触形式存在，本文方法可以无缝兼容现有 NSC 的约束装配与互补求解过程。

\subsection{算法流程}

\noindent\textbf{算法 1：多 Patch / 多 Subpatch 保持 Wrench 等效的 SDF 接触约化}

\noindent 输入：候选接触对 $(A,B)$、SDF 几何、相对状态、误差阈值。\\
\noindent 输出：可注入 Chrono NSC 的 reduced point contacts。

\begin{enumerate}[label=\arabic*.]
    \item 由 SDF 构造高密积分点集合 $\mathcal{Q}$；
    \item 根据连通性与法向一致性得到 patch 集合 $\{\Omega_\ell\}$；
    \item 对每个 patch：
    \begin{enumerate}[label*=\arabic*.]
        \item 初始化 subpatch 队列；
        \item 对每个 subpatch：
        \begin{enumerate}[label*=\arabic*.]
            \item 计算局部平面化误差、二阶矩误差与拓扑误差；
            \item 若误差超阈值，则分裂 subpatch；
            \item 否则初始化代表点 $Q_{\ell r}$；
            \item 优化 $Q_{\ell r}$ 以减小 wrench 锥误差；
            \item 通过局部 dense 微求解得到参考 wrench $w_{\ell r}^\star$；
            \item 解 reduced SOCP，得到代表点力 $\hat{\lambda}$；
            \item 评估 $e_w,e_{\mathcal{C}},e_{\mathrm{CoP}},e_g$；
            \item 若任一误差超阈值，则升阶或细分并返回步骤 3.2.1。
        \end{enumerate}
    \end{enumerate}
    \item 将全部 reduced support points 转换为普通 NSC contact items；
    \item 调用 Chrono NSC 接触容器完成系统装配与互补求解。
\end{enumerate}

\subsection{方法讨论}

本文方法并不试图恢复真实的局部压力场，而是以 dense SDF 接触区域为参考，在\textbf{刚体动力学层面}保持其可行 wrench 锥与当前参考 wrench。相较于基于少量接触点缓存的流形方法，本文方法将接触约化从“缓存少量点”提升为“压缩一个连续接触区域在 wrench 空间中的可行集”；相较于直接采用 patch-level 六维接触元的方法，本文方法又保留了与 Chrono NSC 标准点接触求解链路的兼容性。因此，它特别适用于以下场景：需要高精度 body-level 合力、合力矩与作用线等效，但又希望继续复用现有 NSC 互补求解器的刚体接触动力学问题。

\begin{remark}
本文方法的核心边界在于：它追求的是刚体层面的高精度 wrench 等效，而非局部压力场、接触应力峰值或真实材料变形场的逐点重建。因此，当研究目标转向材料级接触应力与弹性变形主导的接触斑演化时，更适合采用分布式压力场或柔顺接触建模框架。
\end{remark}
